% 原文 p.317
\hypertarget{srcpage:317}{}\label{srcpage:317}
この変換は、いわゆるHadamard回転
\begin{equation}H=\frac1{\sqrt2}\begin{pmatrix}1&1\\1&-1\end{pmatrix}\tag{7.40}\end{equation}
を元の基底状態$|0\rangle$と$|1\rangle$に適用することで行える。基底$\{|0\rangle,|1\rangle\}$の位相反転$\hat\sigma_z$は、基底$\{|+\rangle,|-\rangle\}$ではビット反転誤りに見えることに注意しよう。なぜなら、
\begin{equation}\hat\sigma_z|\pm\rangle=|\mp\rangle\tag{7.41}\end{equation}
だからである。こうして符号化
\begin{equation}\alpha|0\rangle+\beta|1\rangle\longrightarrow\alpha|+++\rangle+\beta|---\rangle\tag{7.42}\end{equation}
を達成した。三つの量子ビットは環境と相互作用し、もつれとデコヒーレンスを受けるため、位相誤り$\hat Z_{ijk}$（$i,j,k\in\{0,1\}$）が生じる。重み1以下の誤り演算子（すなわち$\hat Z_{000}$、$\hat Z_{100}$、$\hat Z_{010}$、$\hat Z_{001}$）だけを考えるとしよう。すると、量子ビットと環境を合わせた状態は
\begin{equation}\begin{aligned}|\Psi\rangle={}&(\alpha|+++\rangle+\beta|---\rangle)|e_{000}\rangle\\&+(\alpha|-++\rangle+\beta|+--\rangle)|e_{100}\rangle\\&+(\alpha|+-+\rangle+\beta|-+-\rangle)|e_{010}\rangle\\&+(\alpha|++-\rangle+\beta|--+\rangle)|e_{001}\rangle\end{aligned}\tag{7.43}\end{equation}
と書ける。以下では、環境状態は書かないことにする。その正確な形には関心がないので、状態中の各量子ビット項が、ある環境の相対状態と相関していると考えればよい。

ここで誤りを診断したい。そのために、まず各量子ビットへHadamard回転(7.40)をもう一度適用し、元の基底$\{|0\rangle,|1\rangle\}$へ戻す。すると状態(7.43)は、（環境状態を省略して）
\begin{equation}\begin{aligned}|\Psi\rangle={}&(\alpha|000\rangle+\beta|111\rangle)\\&+(\alpha|100\rangle+\beta|011\rangle)\\&+(\alpha|010\rangle+\beta|101\rangle)\\&+(\alpha|001\rangle+\beta|110\rangle)\end{aligned}\tag{7.44}\end{equation}
となる。次に、状態$|a_{\mathrm r}\rangle=|00\rangle$に準備した二つの補助量子ビットを導入する。\footnote{実際には、これら二つの補助量子ビットは、誤りの診断と訂正の実行に厳密に必要なわけではないことを示せる。原理的には、誤り症候を、符号化された三つの量子ビット自体から直接取り出すこともできる。しかし、アンシラを含めることで、三ビット符号の説明は、上で概説した量子誤り訂正の基本的な考えを、より明瞭になぞれるようになる。さらに、いわゆる耐故障性をもつ誤り訂正方式を実現するには、補助量子ビットが必要になることがわかる（後の7.4.6節参照）。}量子ビット誤りの情報をアンシラへ符号化するため、次のようにCNOTゲート列を実行する。
% 原文 p.318
\hypertarget{srcpage:318}{}\label{srcpage:318}
\begin{enumerate}
\item まず、Sの第一および第二量子ビットから、第一の補助量子ビットへCNOTゲートを適用する。
\item 次に、Sの第一および第三量子ビットから、第二の補助量子ビットへCNOTゲートを適用する。
\end{enumerate}
三量子ビットの各計算基底状態を入力とした場合に、このCNOTゲート列がアンシラの初期状態$|00\rangle$へ及ぼす効果を、表7.2に示す。上のCNOTゲート列を適用した後、アンシラと三量子ビット系を合わせた状態(7.44)は、
\begin{equation}\begin{aligned}|\Psi\rangle|00\rangle\longrightarrow{}&(\alpha|000\rangle+\beta|111\rangle)|00\rangle\\&+(\alpha|100\rangle+\beta|011\rangle)|11\rangle\\&+(\alpha|010\rangle+\beta|101\rangle)|10\rangle\\&+(\alpha|001\rangle+\beta|110\rangle)|01\rangle\end{aligned}\tag{7.45}\end{equation}
となる。これで、四つの可能性$\hat Z_{000}$、$\hat Z_{100}$、$\hat Z_{010}$、$\hat Z_{001}$のそれぞれを、異なるアンシラ状態$|ij\rangle$（$i,j\in\{0,1\}$）と相関させたことがわかる。これらの状態は互いに直交するので、二つの補助量子ビットの測定によって、四つの誤り型を完全に区別できる。

誤り訂正手続きは、ほぼ完了した。アンシラの測定から、結果（すなわち誤り症候）「01」が得られたとしよう。この測定は、アンシラと量子ビット系を合わせた状態(7.45)を
\begin{equation}(\alpha|001\rangle+\beta|110\rangle)|01\rangle\tag{7.46}\end{equation}
へ収縮させた。アンシラはいまや量子ビット系から完全にもつれが解けたので、以下ではその状態を無視してよい。

測定結果に基づいて、誤りを訂正するために逆変換$\hat I\otimes\hat I\otimes\hat\sigma_x$を適用しなければならないことが、直ちにわかる（表7.3参照）。この変換は状態(7.46)を、次のように変える。
\begin{table}[htbp]\centering\small
\caption{左列に挙げる量子ビットと環境の各入力状態に対して、第一（中央列）および第二（右列）のCNOTゲートを適用した後の、二量子ビットアンシラの初期状態$|00\rangle$の時間発展。}
\begin{tabular}{ccc}\hline 入力&第一のCNOT&第二のCNOT\\\hline
$|000\rangle$&$|00\rangle$&$|00\rangle$\\
$|111\rangle$&$|00\rangle$&$|00\rangle$\\
$|100\rangle$&$|10\rangle$&$|11\rangle$\\
$|011\rangle$&$|10\rangle$&$|11\rangle$\\
$|010\rangle$&$|10\rangle$&$|10\rangle$\\
$|101\rangle$&$|10\rangle$&$|10\rangle$\\
$|001\rangle$&$|00\rangle$&$|01\rangle$\\
$|110\rangle$&$|00\rangle$&$|01\rangle$\\\hline\end{tabular}\end{table}
% 原文 p.319
\hypertarget{srcpage:319}{}\label{srcpage:319}
\begin{table}[htbp]\centering\small
\caption{位相反転誤りを訂正する三ビット符号。二つの補助量子ビットの測定から得られる四つの誤り症候、符号化された三量子ビット状態の対応する誤り、および元の三量子ビット状態の回復に必要な逆変換を示す。}
\begin{tabular}{cll}\hline 症候&誤り&回復変換\\\hline
00&$\hat Z_{000}$（誤りなし）&不要\\
11&$\hat Z_{100}$（第一量子ビットの位相反転）&$\hat\sigma_x\otimes\hat I\otimes\hat I$\\
10&$\hat Z_{010}$（第二量子ビットの位相反転）&$\hat I\otimes\hat\sigma_x\otimes\hat I$\\
01&$\hat Z_{001}$（第三量子ビットの位相反転）&$\hat I\otimes\hat I\otimes\hat\sigma_x$\\\hline\end{tabular}\end{table}
\begin{equation}\alpha|001\rangle+\beta|110\rangle\xrightarrow{\hat I\otimes\hat I\otimes\hat\sigma_x}\alpha|000\rangle+\beta|111\rangle.\tag{7.47}\end{equation}
したがって、最初の冗長に符号化された量子ビット状態(7.39)を回復した。最後に、符号化変換(7.39)の逆を適用すれば、擾乱を受ける前の元の単一量子ビット状態$\alpha|0\rangle+\beta|1\rangle$に到達できる。

二つ以上の量子ビットに影響する位相誤りを含めると（すなわち重み1より大きな位相誤り演算子を考えると）、この位相誤り訂正方式は、つねに成功するわけではない。これは、二つの補助量子ビットから得られる誤り情報が、二つの古典ビットにすぎないことから容易にわかる。したがって、確実に区別できる誤りの可能性は、四つまでである。三量子ビットに可能なすべての位相誤り演算子（重み0から3までのすべての演算子）を含めると、元の状態を回復できるのは、ある確率でのみとなる。今度は八つの異なる誤りを訂正できなければならないため、一つの特定の誤り症候が二つの異なる誤りに対応する。両可能性の確率を計算し、より起こりやすい誤りに対応する逆変換を適用することができる。

単一ビット誤りの確率$p$が50\%未満で、誤りが無相関であると仮定する（二量子ビット誤りと三量子ビット誤りの確率は、それぞれ$3p^2(1-p)$と$p^3$となる）。すると、最も起こりやすい誤り型は、まさに上の例で考えた$\hat Z_{000}$、$\hat Z_{100}$、$\hat Z_{010}$、$\hat Z_{001}$である。より具体的には、この方式の失敗確率は$3p^2(1-p)+p^3$（単に二量子ビット誤りと三量子ビット誤りが起こる確率）なので、失敗確率は$O(p^2)$と変化し、$p$が小さくなるとかなり急速に減少する。これは、まったく訂正しない場合の、失敗確率が$O(p)$と変化する状況と比較される（誤りが一度でも起これば状態は乱され、その確率は定義により$p$である）。したがって上の誤り訂正方式は、$p<50\%$であるかぎり、誤り訂正なしの場合より状況を改善する。この状況は、7.4.1節で論じた古典的誤り訂正の多数決方式に似ている。そこでは、平均して一つを超えるビットが反転しない場合にのみ、誤り回復が利点を与えることを示した。
% 原文 p.320
\hypertarget{srcpage:320}{}\label{srcpage:320}
すなわち、単一ビット反転の確率が50\%未満の場合である。

興味深いことに、ここで述べた位相誤りの三ビット符号は、1998年に実現された、最初の量子誤り訂正の実験的実行\cite{ref361}の基礎となった。\footnote{ただし、この実験は、上で導入した二つの補助量子ビットを用いていない。代わりに、量子ビット系に対する、ある種の共同測定から誤り症候を直接得た。}この実験は、アンサンブル量子計算装置を表すために、液体状態NMRを用いて行われた\cite{ref351}。実験で実行された量子誤り訂正手順は、理論的に予測された量子ビット系からのコヒーレンス喪失を、実際に遅らせることがわかった。

冗長符号化は、符号化しない場合より、量子コンピューター内の物理量子ビットの数を大きく増やす必要がある。そのため、大量のさらなるデコヒーレンスを導入しないよう、非常に注意深く実行しなければならない。（先に、補助量子ビットの文脈でも同様の問題を指摘した。）したがって全体を貫く目標は、物理量子ビットと論理量子ビットの比をできるだけ小さくしながら、誤り訂正能力を最大にすることである。

\subsubsectionlink{装置が引き起こすデコヒーレンスと耐故障性}
これまでは、制御されない環境との相互作用を位相誤りの源と考えてきた。しかし、この文脈の「環境」は、本来、量子ビット系と相互作用するあらゆる物理系という、より広い意味を意図している。その重要な相互作用の一つは、量子ビットを準備し、操作し、測定することで量子情報を処理する装置と、量子ビット系との結合である。装置とのこうした相互作用は、一般に装置と量子ビット系をもつれさせ、そのため量子ビットのデコヒーレンスをもたらす。たとえば、各ゲート操作は通常、ある程度の量子ビットのデコヒーレンスを生じさせる。また、実際のゲート操作はつねに不完全であり、古典的雑音を導入する。量子アルゴリズムには通常、何千ものゲートを量子ビット系へ適用する必要があり、そのため誤りが急速に蓄積する。

したがって、制御されない環境との連続的相互作用による量子ビット誤りだけでなく、量子情報の処理による誤りも周期的に訂正しなければならない。もちろん先に述べたように、誤り訂正自体も装置による測定とゲート操作に基づくため、今度はそれ自体が量子ビットにデコヒーレンス効果を導入し、どのような現実的状況でも無限の精度では実行できない。この問題は、いわゆる耐故障性をもつ誤り検出と回復の方法\cite{ref362,ref363,ref364,ref365,ref366}によって軽減できる。この方法は、個々のゲート操作と測定操作に「雑音」がある（すなわち、ある誤り確率をもつ）という現実的状況でも、量子コンピューターの信頼できる動作を保証しようとする。
